\documentclass[11pt]{article}
\usepackage[utf8]{inputenc}
\usepackage[margin=1in]{geometry}
\usepackage{amsmath}
\usepackage{graphicx}
\usepackage{booktabs}
\usepackage{hyperref}
\usepackage[round]{natbib}
\usepackage{float}

\title{Sensory Coding Transitions from Information Preservation to Discrimination Optimization Along the Visual Dorsal Stream}
\author{Author A$^{1}$, Author B$^{1,2}$, Author C$^{2}$ \\
\small $^{1}$Department of Neuroscience, University \\
\small $^{2}$Center for Neural Science, University}
\date{}

\begin{document}
\maketitle

\begin{abstract}
Optimal coding theory predicts that early sensory areas maximize information about stimulus distributions, but how this optimization changes along the cortical processing hierarchy is poorly understood. Here we re-analyzed 1,056 binocular disparity tuning curves from macaque V1 ($n \approx 400$), V2 ($n \approx 300$), and MT ($n \approx 350$), combined with natural binocular disparity statistics from eye-tracked humans performing food preparation and navigation tasks. We found that the power law exponent $\alpha$ relating population Fisher information (FI) to stimulus probability shifts from approximately 2 (information maximization, ``infomax'') in V1 and V2 to approximately 2/3 (discrimination optimization, ``discrimax'') in MT. V1 and V2 exponent distributions were highly overlapping (median $\alpha \approx 2.0$), while MT was significantly shifted ($\alpha \approx 0.67$; $p < 0.001$ vs. V1 and V2). Fitting each tuning curve with a 6-parameter modified Gabor function and systematically swapping individual parameters between areas revealed that the preferred disparity parameter (Gaussian envelope center) is the primary driver of the coding transition: replacing V1 preferred disparities with MT values shifted the V1 population FI toward the discrimination-optimal regime. These results demonstrate a progressive transformation of sensory coding strategy along the dorsal visual stream, from preserving all information about the stimulus distribution to optimizing perceptual discrimination performance.
\end{abstract}

\section{Introduction}

A central principle of sensory neuroscience is that neural populations encode information about stimuli in a manner shaped by the statistical structure of the natural environment \citep{barlow1961coding, attneave1954informational, simoncelli2001natural}. The efficient coding hypothesis predicts that neural resources should be allocated to maximize information transmission, with more neurons and finer coding resolution devoted to statistically more probable stimuli \citep{laughlin1981simple}.

Formally, this relationship can be quantified through Fisher information (FI), which measures the precision with which a neural population encodes a stimulus \citep{fisher_information1925}. For a population that maximizes information transmission (``infomax'' coding), the relationship between population FI and stimulus probability $p(s)$ follows a power law: $\text{FI}(s) \propto p(s)^{2}$ \citep{ganguli2012efficient, wei2015bayesian}. This exponent of 2 corresponds to minimizing the $L_0$ error norm, which penalizes all estimation errors equally regardless of magnitude.

However, perception ultimately serves behavior, and behavioral performance depends not on information preservation per se but on the ability to discriminate between stimuli \citep{series2009tuning}. A discrimination-optimal code follows a different power law: $\text{FI}(s) \propto p(s)^{2/3}$ \citep{ganguli2012efficient}, corresponding to minimizing the $L_2$ error norm (mean squared error), which penalizes large estimation errors disproportionately. Whether and where the brain transitions from infomax to discrimination-optimal coding along the cortical hierarchy has not been tested empirically.

Here we address this question using binocular disparity---the horizontal difference between left and right retinal images---as a model sensory dimension. Disparity is encoded by neurons throughout the dorsal visual stream, with well-characterized tuning properties in V1, V2, and MT \citep{cumming2001binocular, deangelis1998cortical, prince2002quantitative}. By combining a large dataset of disparity tuning curves with natural disparity statistics, we test whether population coding strategy shifts from infomax to discrimax along the V1--V2--MT hierarchy.

\section{Results}

\subsection{Natural Binocular Disparity Statistics}

Natural disparity distributions were computed from eye-tracked humans ($n = 3$ adults) performing food preparation and navigation tasks within 10 degrees of fixation (Table~\ref{tab:disparity_stats}).

\begin{table}[H]
\centering
\caption{Properties of natural binocular disparity distributions by task. Distribution parameters computed from 100 bootstrapped resamples \citep{bootstrap_efron1979}.}
\label{tab:disparity_stats}
\begin{tabular}{lcc}
\toprule
Property & Food Preparation & Navigation \\
\midrule
Mean (arcmin) & $-$0.12 $\pm$ 0.08 & 0.05 $\pm$ 0.06 \\
Symmetry & Approximately symmetric & Approximately symmetric \\
Kurtosis & 8.42 $\pm$ 0.94 & 12.18 $\pm$ 1.31 \\
\% within $\pm$30 arcmin & 82.4 & 91.7 \\
Large disparities & More frequent & Less frequent \\
\bottomrule
\end{tabular}
\end{table}

Both distributions were approximately zero-mean and symmetric, with high kurtosis indicating that small disparities are much more probable than large disparities. The food preparation task generated more large disparities (close objects being manipulated) than navigation (distant scenes).

\subsection{Population Fisher Information Across Brain Areas}

Population FI was computed from 1,056 tuning curves across V1, V2, and MT, assuming Poisson-like noise \citep{poisson_noise2009} (Table~\ref{tab:neuron_counts}).

\begin{table}[H]
\centering
\caption{Neural dataset: tuning curve counts and recording conditions.}
\label{tab:neuron_counts}
\begin{tabular}{lccc}
\toprule
Brain Area & Neuron Count & Source Studies & Recording Condition \\
\midrule
V1 & $\sim$400 & Multiple & Awake fixating macaque \\
V2 & $\sim$300 & Multiple & Awake fixating macaque \\
MT & $\sim$350 & Multiple & Awake fixating macaque \\
Total & 1,056 & -- & -- \\
\bottomrule
\end{tabular}
\end{table}

Population FI distributions showed a progressive change across areas: V1 and V2 FI was strongly concentrated near zero disparity (tracking the high probability of small disparities), while MT FI was more spread and flatter (Table~\ref{tab:fi_shape}).

\begin{table}[H]
\centering
\caption{Population FI distribution shape metrics (500 bootstrap samples per area). FI concentration ratio = FI at 0 arcmin / FI at 60 arcmin.}
\label{tab:fi_shape}
\begin{tabular}{lccc}
\toprule
Metric & V1 & V2 & MT \\
\midrule
FI concentration ratio & 14.2 $\pm$ 2.1 & 12.8 $\pm$ 1.9 & $\mathbf{4.3 \pm 0.8}$ \\
FI kurtosis & 6.8 $\pm$ 0.7 & 6.2 $\pm$ 0.6 & $\mathbf{2.9 \pm 0.4}$ \\
FI peak width (FWHM, arcmin) & 22.4 $\pm$ 1.8 & 24.1 $\pm$ 2.0 & $\mathbf{48.7 \pm 3.4}$ \\
\bottomrule
\end{tabular}
\end{table}

\subsection{Power Law Exponent Shifts from V1/V2 to MT}

The key test was the power law exponent $\alpha$ in the relationship $\text{FI}(s) \propto p(s)^{\alpha}$ (Table~\ref{tab:alpha}).

\begin{table}[H]
\centering
\caption{Power law exponent $\alpha$ distributions by brain area and task (500 bootstrapped samples per condition, Gaussian fits). Theoretical predictions: infomax $\alpha = 2$; discrimax $\alpha = 2/3$.}
\label{tab:alpha}
\begin{tabular}{lcccccc}
\toprule
 & \multicolumn{3}{c}{Food Preparation} & \multicolumn{3}{c}{Navigation} \\
\cmidrule(lr){2-4} \cmidrule(lr){5-7}
Area & Median $\alpha$ & SD & 95\% CI & Median $\alpha$ & SD & 95\% CI \\
\midrule
V1 & 1.94 & 0.28 & [1.41, 2.48] & 2.08 & 0.31 & [1.49, 2.68] \\
V2 & 1.87 & 0.26 & [1.38, 2.37] & 1.96 & 0.29 & [1.41, 2.52] \\
MT & $\mathbf{0.71}$ & 0.19 & [0.35, 1.08] & $\mathbf{0.64}$ & 0.21 & [0.24, 1.04] \\
\bottomrule
\end{tabular}
\end{table}

V1 and V2 exponent distributions were centered near the infomax prediction ($\alpha = 2$) and highly overlapping. MT exponents were dramatically shifted toward the discrimax prediction ($\alpha = 2/3$). Statistical comparisons confirmed that V1 versus V2 was non-significant (overlapping 95\% CIs), while V1 versus MT and V2 versus MT were highly significant ($p < 0.001$ by bootstrap permutation test).

\subsection{Gabor Parameter Analysis Identifies the Driver of Coding Transition}

Each tuning curve was fit with a 6-parameter modified Gabor function: Gaussian center ($\mu$, preferred disparity), width ($\sigma$), amplitude ($A$), cosine frequency ($f$), phase ($\phi$), and baseline ($b$) (Table~\ref{tab:gabor}).

\begin{table}[H]
\centering
\caption{Gabor function parameter distributions by brain area. Median [IQR]. Kruskal-Wallis test across areas.}
\label{tab:gabor}
\begin{tabular}{lcccc}
\toprule
Parameter & V1 & V2 & MT & $p$ (K-W) \\
\midrule
$\mu$ (arcmin) & 2.1 [0.8, 5.4] & 3.2 [1.1, 7.8] & $\mathbf{12.4}$ [4.8, 28.6] & $< 0.001$ \\
$\sigma$ (arcmin) & 8.3 [4.2, 14.1] & 10.7 [5.1, 18.3] & $\mathbf{22.8}$ [11.4, 38.2] & $< 0.001$ \\
$f$ (cycles/arcmin) & 0.042 [0.021, 0.078] & 0.034 [0.018, 0.062] & $\mathbf{0.018}$ [0.008, 0.034] & $< 0.001$ \\
$\phi$ (rad) & 1.24 [0.42, 2.18] & 1.31 [0.48, 2.24] & 1.18 [0.38, 2.12] & 0.62 \\
$A$ (spikes/s) & 18.4 [8.2, 32.1] & 22.1 [9.8, 38.4] & 24.8 [12.1, 42.6] & 0.14 \\
$b$ (spikes/s) & 6.2 [2.1, 12.4] & 7.8 [2.8, 14.2] & 8.4 [3.2, 15.8] & 0.28 \\
\bottomrule
\end{tabular}
\end{table}

MT neurons showed significantly larger preferred disparities and broader tuning widths than V1/V2, while phase, amplitude, and baseline did not differ significantly across areas.

\subsection{Parameter Swapping Reveals Preferred Disparity as Key Driver}

To identify which parameter drives the V1-to-MT coding transition, we systematically replaced each V1 parameter with randomly sampled MT values and recomputed population FI (Table~\ref{tab:swap}).

\begin{table}[H]
\centering
\caption{Parameter swapping analysis: effect of replacing individual V1 parameters with MT values on the population FI power law exponent. 500 resampled V1 populations per swap. $\Delta\alpha$ = change from original V1 $\alpha$ toward MT $\alpha$.}
\label{tab:swap}
\begin{tabular}{lcccc}
\toprule
Parameter Swapped & Resulting $\alpha$ & $\Delta\alpha$ & \% of V1$\to$MT shift & Key Driver? \\
\midrule
None (original V1) & 1.94 & 0.00 & 0 & -- \\
$\mu$ (preferred disparity) & $\mathbf{0.92}$ & $\mathbf{-1.02}$ & $\mathbf{83\%}$ & Yes \\
$\sigma$ (tuning width) & 1.61 & $-$0.33 & 27\% & Partial \\
$f$ (frequency) & 1.82 & $-$0.12 & 10\% & No \\
$\phi$ (phase) & 1.91 & $-$0.03 & 2\% & No \\
$A$ (amplitude) & 1.90 & $-$0.04 & 3\% & No \\
$b$ (baseline) & 1.89 & $-$0.05 & 4\% & No \\
MT original & 0.71 & $-$1.23 & 100\% & -- \\
\bottomrule
\end{tabular}
\end{table}

Replacing V1 preferred disparities with MT values accounted for 83\% of the coding transition, shifting $\alpha$ from 1.94 to 0.92. Tuning width contributed a partial additional shift (27\%), while the remaining parameters had minimal effects ($<$10\% each). This identifies the redistribution of preferred disparities toward larger values as the primary mechanism underlying the infomax-to-discrimax transition.

\section{Discussion}

Our analysis reveals a progressive transformation of sensory coding strategy along the visual dorsal stream: V1 and V2 populations encode binocular disparity in an information-maximizing regime ($\alpha \approx 2$), while MT shifts to a discrimination-optimal regime ($\alpha \approx 2/3$). This shift is consistent with theoretical predictions from efficient coding theory \citep{ganguli2012efficient, wei2015bayesian} and demonstrates that the brain implements different optimization objectives at different hierarchical levels.

The infomax strategy in V1/V2 is well-suited for early visual processing, where the goal is to preserve as much information as possible about the full stimulus distribution. By allocating FI proportional to $p(s)^2$, these populations dedicate the most coding resources to the most frequently encountered disparities (near zero), minimizing information loss \citep{laughlin1981simple, barlow1961coding}.

The discrimax strategy in MT, by contrast, allocates FI more evenly across the disparity range (proportional to $p(s)^{2/3}$), which optimizes the ability to discriminate between stimuli by minimizing mean squared estimation error \citep{ganguli2012efficient}. This is functionally appropriate for MT, which is critically involved in perceptual judgments about depth and 3D motion \citep{deangelis1998cortical, czuba2014area}.

The parameter swapping analysis identifies a simple mechanism for this transition: MT neurons prefer larger disparities than V1/V2 neurons. This shift in preferred disparity away from zero spreads the population's sensitivity across a wider range, naturally transitioning the FI--probability relationship from infomax to discrimax without requiring changes in other tuning parameters. The broader tuning widths in MT contribute partially but are not the primary driver.

Several limitations should be noted. The independence assumption (no noise correlations) may not hold \citep{haefner2013inferring}, though this is unavoidable with single-unit data compiled across studies. Resource constraints are assumed similar across areas, which may not be strictly true given differences in cell density and metabolic demands. The analysis is restricted to binocular disparity, and whether similar coding transitions occur for other stimulus dimensions remains to be tested.

\section{Methods}

\subsection{Natural Disparity Statistics}

Binocular disparity distributions were computed from a previously collected dataset of 3 adults wearing binocular eye-tracking glasses during food preparation and navigation tasks. Probability distributions were computed within 10 degrees of fixation and resampled based on receptive field center distributions of each neural population to account for hemifield recording biases.

\subsection{Neural Data}

Disparity tuning curves from 1,056 isolated single units in awake fixating macaque monkeys were compiled from published datasets spanning V1, V2, and MT. Each tuning curve was fit with a 6-parameter modified Gabor function \citep{gabor_function1946}: $r(d) = b + A \cdot \exp(-(d - \mu)^2 / 2\sigma^2) \cdot \cos(2\pi f(d - \mu) + \phi)$.

\subsection{Fisher Information Computation}

For each neuron, individual FI was computed assuming Poisson-like variability: $\text{FI}_i(s) = [r'_i(s)]^2 / r_i(s)$, where $r_i(s)$ is the mean firing rate and $r'_i(s)$ is its derivative. Population FI was the sum of individual FI across the sampled population. The power law exponent $\alpha$ was estimated by log-log regression of population FI against disparity probability. Bootstrapping (500 samples) generated distributions of $\alpha$.

\subsection{Parameter Swapping}

For each parameter, 500 resampled V1 populations were generated by replacing that parameter's values with random draws from the MT distribution while keeping all other parameters at their V1 values. Population FI and the resulting $\alpha$ were computed for each resampled population.

\section{Conclusion}

We demonstrate that the coding strategy for binocular disparity transitions from information maximization in V1/V2 to discrimination optimization in MT, driven primarily by a shift in preferred disparities toward larger values. This finding provides empirical support for theoretical predictions about hierarchical coding transformations and reveals that the visual dorsal stream progressively reshapes sensory representations from information preservation toward perceptual utility.

\bibliographystyle{plainnat}
\bibliography{references}

\end{document}
